\documentclass[UTF8,fontset=fandol,11pt]{ctexart}
\usepackage[a4paper,margin=24mm]{geometry}
\usepackage{amsmath,amssymb,amsthm,graphicx,tabularx,booktabs,hyperref}
\newtheorem{definition}{定义}[section]
\newtheorem{theorem}[definition]{定理}
\newtheorem{example}[definition]{例题}
\newtheorem{exercise}[definition]{练习}
\hypersetup{unicode=true,colorlinks=true,linkcolor=blue}
\title{函数极限与连续：本科一年级微积分讲义计划}
\date{}
\begin{document}
\maketitle
\section{一元函数极限：定义、左右极限与无穷远极限}\label{sec:S01}
% UNIT S01.U01
\begin{definition}
设函数 $f$ 在点 $x_0$ 的某个去心邻域内有定义。若存在常数 $A$，使得对任意 $\varepsilon>0$，都存在 $\delta>0$，满足当 $0<|x-x_0|<\delta$ 时有 $|f(x)-A|<\varepsilon$，则称 $A$ 为 $f(x)$ 当 $x\to x_0$ 时的极限，记为
\[
\lim_{x\to x_0}f(x)=A.
\]
其中，$x_0$ 不要求属于函数的定义域。\label{def:point-limit}
\end{definition}
% UNIT S01.U02
\begin{definition}
设函数 $f$ 在 $x_0$ 的某个左侧去心邻域内有定义。若存在常数 $A_-$，使得对任意 $\varepsilon>0$，都存在 $\delta>0$，当 $x_0-\delta<x<x_0$ 时有 $|f(x)-A_-|<\varepsilon$，则称 $A_-$ 为 $f(x)$ 当 $x\to x_0^-$ 时的左极限，记为
\[
\lim_{x\to x_0^-}f(x)=A_-.
\]
设函数 $f$ 在 $x_0$ 的某个右侧去心邻域内有定义。若存在常数 $A_+$，使得对任意 $\varepsilon>0$，都存在 $\delta>0$，当 $x_0<x<x_0+\delta$ 时有 $|f(x)-A_+|<\varepsilon$，则称 $A_+$ 为 $f(x)$ 当 $x\to x_0^+$ 时的右极限，记为
\[
\lim_{x\to x_0^+}f(x)=A_+.
\]
\label{def:one-sided-limit}
\end{definition}
% UNIT S01.U03
\begin{definition}
设函数 $f$ 在某个右半无限区间 $(M,+\infty)$ 上有定义。若存在常数 $A$，使得对任意 $\varepsilon>0$，都存在 $R>M$，当 $x>R$ 时有 $|f(x)-A|<\varepsilon$，则称 $A$ 为 $f(x)$ 当 $x\to+\infty$ 时的极限，记为
\[
\lim_{x\to+\infty}f(x)=A.
\]
类似地，若函数在某个左半无限区间 $(-\infty,M)$ 上有定义，且对任意 $\varepsilon>0$ 都存在 $R<M$，使得当 $x<R$ 时有 $|f(x)-A|<\varepsilon$，则记为
\[
\lim_{x\to-\infty}f(x)=A.
\]
\label{def:infinite-limit}
\end{definition}
% UNIT S01.U04
\begin{theorem}
设函数 $f$ 在点 $x_0$ 的某个去心邻域内有定义，则
\[
\lim_{x\to x_0}f(x)=A
\]
当且仅当
\[
\lim_{x\to x_0^-}f(x)=A
\quad\text{且}\quad
\lim_{x\to x_0^+}f(x)=A.
\]
也就是说，函数在一点的双侧极限存在，当且仅当该点的左极限和右极限都存在且相等。\label{thm:two-sided-limit-criterion}
\end{theorem}
% UNIT S01.U05
\begin{proof}
先设 $\lim_{x\to x_0}f(x)=A$。给定任意 $\varepsilon>0$，由双侧极限的定义，存在 $\delta>0$，使得当 $0<|x-x_0|<\delta$ 时有 $|f(x)-A|<\varepsilon$。特别地，当 $x_0-\delta<x<x_0$ 或 $x_0<x<x_0+\delta$ 时，上述不等式分别成立，因此左极限和右极限都等于 $A$。

反过来，设左极限和右极限都等于 $A$。给定任意 $\varepsilon>0$，存在 $\delta_1>0$，使得当 $x_0-\delta_1<x<x_0$ 时有 $|f(x)-A|<\varepsilon$；存在 $\delta_2>0$，使得当 $x_0<x<x_0+\delta_2$ 时有 $|f(x)-A|<\varepsilon$。取 $\delta=\min\{\delta_1,\delta_2\}$，则当 $0<|x-x_0|<\delta$ 时，$x$ 必在 $x_0$ 的左侧或右侧，从而都有 $|f(x)-A|<\varepsilon$。因此 $\lim_{x\to x_0}f(x)=A$。
\end{proof}
% UNIT S01.U06
\begin{theorem}
设函数 $f$ 在点 $x_0$ 的某个去心邻域内有定义。若
\[
\lim_{x\to x_0}f(x)=A
\quad\text{且}\quad
\lim_{x\to x_0}f(x)=B,
\]
则 $A=B$。因此，函数在同一点处的有限极限至多有一个。\label{thm:limit-uniqueness}
\end{theorem}
% UNIT S01.U07
\begin{proof}
假设 $A\ne B$。取 $\varepsilon=|A-B|/3>0$。由 $\lim_{x\to x_0}f(x)=A$，存在 $\delta_1>0$，使得当 $0<|x-x_0|<\delta_1$ 时有 $|f(x)-A|<\varepsilon$；由 $\lim_{x\to x_0}f(x)=B$，存在 $\delta_2>0$，使得当 $0<|x-x_0|<\delta_2$ 时有 $|f(x)-B|<\varepsilon$。

取 $\delta=\min\{\delta_1,\delta_2\}$，并取一个满足 $0<|x-x_0|<\delta$ 的 $x$。由三角不等式，
\[
|A-B|\leq |A-f(x)|+|f(x)-B|<2\varepsilon=\frac{2}{3}|A-B|,
\]
这与 $|A-B|>0$ 矛盾。因此 $A=B$。
\end{proof}
% UNIT S01.U08
\begin{example}
计算
\[
\lim_{x\to2}(3x-1)\text{。}
\]
为说明双侧极限的处理方法，分别考察 $x\to2^-$ 和 $x\to2^+$ 两种情形。
\end{example}
% UNIT S01.U09
\par\noindent\textbf{解：}当 $x\to2^-$ 或 $x\to2^+$ 时，都有 $3x-1\to3\cdot2-1=5$。更直接地，对任意 $\varepsilon>0$，取 $\delta=\varepsilon/3$。当 $0<|x-2|<\delta$ 时，
\[
|(3x-1)-5|=|3x-6|=3|x-2|<3\delta=\varepsilon\text{。}
\]
因此由极限定义，
\[
\lim_{x\to2}(3x-1)=5\text{。}
\]\par
% UNIT S01.U10
\begin{example}
设
\[
f(x)=
\begin{cases}
1,&x<0,\\
2,&x>0.
\end{cases}
\]
判断 $\lim_{x\to0}f(x)$ 是否存在。
\end{example}
% UNIT S01.U11
\par\noindent\textbf{解：}当 $x\to0^-$ 时，始终有 $f(x)=1$，所以
\[
\lim_{x\to0^-}f(x)=1\text{。}
\]
当 $x\to0^+$ 时，始终有 $f(x)=2$，所以
\[
\lim_{x\to0^+}f(x)=2\text{。}
\]
由于左右极限不相等，根据定理 \ref{thm:two-sided-limit-criterion}，$\lim_{x\to0}f(x)$ 不存在。\par
% UNIT S01.U12
\begin{exercise}
\label{mathagent:exercise:bf60311c135cf0c14f3cb42069a931ed175758be7d89dff861bed83a8d332803}
设
\[
f(x)=
\begin{cases}
2x+1,&x<1,\\
5,&x>1.
\end{cases}
\]
判断 $\lim_{x\to1}f(x)$ 是否存在；若存在，求出该极限。
\end{exercise}
\section*{练习答案}
% UNIT S01.U13
\par\noindent\textbf{\hyperref[mathagent:exercise:bf60311c135cf0c14f3cb42069a931ed175758be7d89dff861bed83a8d332803]{练习 \ref*{mathagent:exercise:bf60311c135cf0c14f3cb42069a931ed175758be7d89dff861bed83a8d332803} 的解答：}}当 $x\to1^-$ 时，$f(x)=2x+1$，因此
\[
\lim_{x\to1^-}f(x)=2\cdot1+1=3\text{。}
\]
当 $x\to1^+$ 时，$f(x)=5$，因此
\[
\lim_{x\to1^+}f(x)=5\text{。}
\]
由于左极限 $3$ 与右极限 $5$ 不相等，根据定理 \ref{thm:two-sided-limit-criterion}，双侧极限 $\lim_{x\to1}f(x)$ 不存在。\par
\end{document}
